% GENERATED FILE. Edit canonical lessons, then run npm run generate:books.
% canonical-source-hash: fnv1a64:79f09e452b715e15
% canonical-lessons: TA-C290-kopam-vantatu, TA-C290-acai, TA-C290-ilakku, TA-C290-poka-acai, TA-R290-feelings-and-hopes

\chapter{Feelings and Hopes}
\label{ch:ta-feelings-and-hopes}
\hlchaptermodality{290}

\begin{chapteropening}
\textbf{By the end of this chapter:} \emph{I can say how I felt with eṉakkuk kōpam vantatu and eṉakkuc cantōṣamāka iruntatu, say what I long for with ācai and virumpukiṟēṉ, and name a goal (ilakku) and what I want to become.}

Complete the last lesson of chapter 290: say how you felt, what you long for and what you aim to become.
\end{chapteropening}

\section[\texorpdfstring{\ta{எனக்குக்} \ta{கோபம்} \ta{வந்தது} (eṉakkuk kōpam …)}{எனக்குக் கோபம் வந்தது (eṉakkuk kōpam …)}]{\ta{எனக்குக்} \ta{கோபம்} \ta{வந்தது} (eṉakkuk kōpam vantatu) — how I felt}

\label{lesson:TA-C290-kopam-vantatu}



\begin{quote}
Say the words for anger, fear and gladness. (\emph{kōpam}, \emph{payam}, \emph{cantōṣam}.) Then \textquotedblleft{}I know Tamil\textquotedblright{}. (\emph{eṉakkut tamiḻ teriyum}.) Then say \textquotedblleft{}dampness\textquotedblright{}. (\emph{īram}.)
\end{quote}

\subsection*{You'll want to know: a feeling comes to you}
\begin{quote}
\textbf{\ta{எனக்குக்} \ta{கோபம்} \ta{வந்தது}.}
\end{quote}

\begin{quote}
\emph{eṉakkuk kōpam vantatu.} — \textquotedblleft{}I got angry.\textquotedblright{}
\end{quote}

Word for word: \textquotedblleft{}to me, anger came\textquotedblright{}. The feeling is what came, so the verb is the \textquotedblleft{}it\textquotedblright{} past, \textbf{\ta{வந்தது}}, and the person it came to takes \textbf{-\ta{க்கு}}, as in \textbf{\ta{எனக்குத்} \ta{தமிழ்} \ta{தெரியும்}}. The first letter of the feeling doubles in writing: \textbf{\ta{எனக்குக்} \ta{கோபம்}}, \textbf{\ta{அவனுக்குப்} \ta{பயம்}}.

\begin{quote}
\textbf{\ta{அவனுக்குப்} \ta{பயம்} \ta{வந்தது}.}
\end{quote}

\begin{quote}
\emph{avaṉukkup payam vantatu.} — \textquotedblleft{}He got frightened.\textquotedblright{}
\end{quote}

\textbf{\ta{அவனுக்கு}} (\emph{avaṉukku}, \textquotedblleft{}to him\textquotedblright{}), \textbf{\ta{அவளுக்கு}} (\emph{avaḷukku}, \textquotedblleft{}to her\textquotedblright{}), \textbf{\ta{உங்களுக்கு}} (\emph{uṅkaḷukku}, \textquotedblleft{}to you\textquotedblright{}).

\subsection*{You'll want to know: the state you are in, -\ta{ஆக} \ta{இருக்கிறது}}
\begin{quote}
\textbf{\ta{எனக்குச்} \ta{சந்தோஷமாக} \ta{இருக்கிறது}.}
\end{quote}

\begin{quote}
\emph{eṉakkuc cantōṣamāka irukkiṟatu.} — \textquotedblleft{}I am happy.\textquotedblright{}
\end{quote}

For a feeling you are in, the feeling takes \textbf{-\ta{ஆக}}, the \textbf{-\ta{ஆக}} of \textbf{\ta{நன்றாக}}, and the verb is \textbf{\ta{இருக்கிறது}}: \textquotedblleft{}to me, it is in gladness\textquotedblright{}. In the past: \textbf{\ta{இருந்தது}}.

\begin{quote}
\textbf{\ta{எனக்குப்} \ta{பயமாக} \ta{இருந்தது}.}
\end{quote}

\begin{quote}
\emph{eṉakkup payamāka iruntatu.} — \textquotedblleft{}I was afraid.\textquotedblright{}
\end{quote}

\begin{quote}
\textbf{\ta{அவளுக்கு} \ta{ஏமாற்றமாக} \ta{இருந்தது}.}
\end{quote}

\begin{quote}
\emph{avaḷukku ēmāṟṟamāka iruntatu.} — \textquotedblleft{}She was disappointed.\textquotedblright{}
\end{quote}

The same feelings also go with a plain subject: \textbf{\ta{நான்} \ta{சந்தோஷமாக} \ta{இருக்கிறேன்}} (\emph{nāṉ cantōṣamāka irukkiṟēṉ}, \textquotedblleft{}I am happy\textquotedblright{}).

\subsection*{Your turn}
Say these aloud:

\begin{itemize}
\raggedright
  \item \textquotedblleft{}I got angry\textquotedblright{} — \emph{eṉakkuk kōpam vantatu}
  \item \textquotedblleft{}He got frightened\textquotedblright{} — \emph{avaṉukkup payam vantatu}
  \item \textquotedblleft{}I am happy\textquotedblright{} — \emph{eṉakkuc cantōṣamāka irukkiṟatu}
  \item \textquotedblleft{}She was disappointed\textquotedblright{} — \emph{avaḷukku ēmāṟṟamāka iruntatu}
\end{itemize}

\begin{tcolorbox}[breakable,colback=teal!4,colframe=teal!35!black,title={Before you move on}]
\textquotedblleft{}I got angry\textquotedblright{}? (\textbf{\ta{எனக்குக்} \ta{கோபம்} \ta{வந்தது}}.) Who takes \textbf{-\ta{க்கு}}? (\textbf{The person who feels.}) \textquotedblleft{}I was afraid\textquotedblright{}? (\textbf{\ta{எனக்குப்} \ta{பயமாக} \ta{இருந்தது}}.) Which verb says the feeling arrived, and which says you were in it? (\textbf{\ta{வந்தது}} arrived; \textbf{\ta{இருந்தது}} was.)
\end{tcolorbox}
\section[\texorpdfstring{\ta{ஆசை} (ācai)}{ஆசை (ācai)}]{\ta{ஆசை} (ācai) — a wish, a longing}

\label{lesson:TA-C290-acai}



\begin{quote}
Before the new one: say \textquotedblleft{}I am happy.\textquotedblright{} (\emph{eṉakkuc cantōṣamāka irukkiṟatu}.) Then \textquotedblleft{}to dream\textquotedblright{} and \textquotedblleft{}the sea\textquotedblright{}. (\emph{kaṉavu kāṇ}, \emph{kaṭal}.)
\end{quote}

\subsection*{You'll want to know: \ta{ஆசை}}
\textbf{\ta{ஆசை}} — \emph{ācai} — \textquotedblleft{}a wish, a longing, a desire\textquotedblright{}.

\begin{quote}
\textbf{\ta{எனக்கு} \ta{ஒரு} \ta{ஆசை} \ta{இருக்கிறது}.}
\end{quote}

\begin{quote}
\emph{eṉakku oru ācai irukkiṟatu.} — \textquotedblleft{}I have a wish.\textquotedblright{}
\end{quote}

\begin{quote}
\textbf{\ta{எனக்குக்} \ta{கடலில்} \ta{நீந்த} \ta{ஆசை}.}
\end{quote}

\begin{quote}
\emph{eṉakkuk kaṭalil nīnta ācai.} — \textquotedblleft{}I long to swim in the sea.\textquotedblright{}
\end{quote}

\textbf{\ta{ஆசை}} goes the way of the feelings: the one who wishes takes \textbf{-\ta{க்கு}}. What you long to do comes as a verb's to-form before it: \textbf{\ta{நீந்த} \ta{ஆசை}}, \textquotedblleft{}a longing to swim\textquotedblright{}. It covers a child's wish for a sweet and a grown person's hope for a whole life; the sentence tells you which.

\begin{cousinweb}
\textbf{\ta{ஆசை}} (\emph{ācai}) is a loan from Sanskrit \emph{āśā}, which Monier-Williams glosses \textquotedblleft{}wish, desire, hope\textquotedblright{}. The final \emph{-ā} of the Sanskrit word became \textbf{-\ta{ஐ}} in Tamil, the change you can hear in \textbf{\ta{கதை}}, \textquotedblleft{}a story\textquotedblright{}, from Sanskrit \emph{kathā}.
\end{cousinweb}

\subsection*{Your turn}
Say these aloud:

\begin{itemize}
\raggedright
  \item \emph{ācai}, slowly, then at speed
  \item \textquotedblleft{}I have a wish\textquotedblright{} — \emph{eṉakku oru ācai irukkiṟatu}
  \item \textquotedblleft{}I long to swim in the sea\textquotedblright{} — \emph{eṉakkuk kaṭalil nīnta ācai}
\end{itemize}

\begin{tcolorbox}[breakable,colback=teal!4,colframe=teal!35!black,title={Before you move on}]
What does \textbf{\ta{ஆசை}} mean? (\textbf{A wish, a longing.}) \textquotedblleft{}I long to swim in the sea\textquotedblright{}? (\textbf{\ta{எனக்குக்} \ta{கடலில்} \ta{நீந்த} \ta{ஆசை}}.) Which language is it from? (\textbf{Sanskrit}: \emph{āśā}.)
\end{tcolorbox}
\section[\texorpdfstring{\ta{இலக்கு} (ilakku)}{இலக்கு (ilakku)}]{\ta{இலக்கு} (ilakku) — a goal, a target}

\label{lesson:TA-C290-ilakku}



\begin{quote}
Before the new one: say \textquotedblleft{}I have a wish.\textquotedblright{} (\emph{eṉakku oru ācai irukkiṟatu}.) Then the words for a doctor and a teacher. (\emph{maruttuvar}, \emph{āciriyar}.)
\end{quote}

\subsection*{You'll want to know: \ta{இலக்கு}}
\textbf{\ta{இலக்கு}} — \emph{ilakku} — \textquotedblleft{}a goal, a target\textquotedblright{}.

\begin{quote}
\textbf{\ta{என்} \ta{இலக்கு} \ta{மருத்துவர்} \ta{ஆவது}.}
\end{quote}

\begin{quote}
\emph{eṉ ilakku maruttuvar āvatu.} — \textquotedblleft{}My goal is to become a doctor.\textquotedblright{}
\end{quote}

\textbf{\ta{ஆவது}} (\emph{āvatu}, \textquotedblleft{}becoming\textquotedblright{}). Where \textbf{\ta{ஆசை}} is a longing, \textbf{\ta{இலக்கு}} is a mark you aim at and work towards. The same word names the target an arrow is shot at.

\begin{quote}
\textbf{\ta{இந்த} \ta{ஆண்டு} \ta{என்} \ta{இலக்கு}: \ta{மாதம்} \ta{ஒரு} \ta{புத்தகம்}.}
\end{quote}

\begin{quote}
\emph{inta āṇṭu eṉ ilakku: mātam oru puttakam.} — \textquotedblleft{}My goal this year: a book a month.\textquotedblright{}
\end{quote}

\textbf{\ta{இந்த}} (\emph{inta}, \textquotedblleft{}this\textquotedblright{}): the \textbf{i-} of \textbf{\ta{இது}} and \textbf{\ta{இங்கே}}, in front of a noun.

\begin{cousinweb}
\textbf{\ta{இலக்கு}} (\emph{ilakku}) is usually traced to Sanskrit \emph{lakṣya} or \emph{lakṣa}, \textquotedblleft{}mark, aim\textquotedblright{}: CDIAL 10881 lists \emph{lakṣá} \textquotedblleft{}mark, sign\textquotedblright{} and \textquotedblleft{}aim\textquotedblright{}, and \emph{lakṣya-} \textquotedblleft{}aim\textquotedblright{}. The \textbf{\ta{இ}} in front is the vowel older Tamil set before a borrowed word that began with \emph{l}.
\end{cousinweb}

\subsection*{Your turn}
Say these aloud:

\begin{itemize}
\raggedright
  \item \emph{ilakku}, slowly, then at speed
  \item \textquotedblleft{}My goal is to become a doctor\textquotedblright{} — \emph{eṉ ilakku maruttuvar āvatu}
  \item \textquotedblleft{}My goal this year: a book a month\textquotedblright{} — \emph{inta āṇṭu eṉ ilakku: mātam oru puttakam}
\end{itemize}

\begin{tcolorbox}[breakable,colback=teal!4,colframe=teal!35!black,title={Before you move on}]
What does \textbf{\ta{இலக்கு}} mean? (\textbf{A goal, a target.}) And \textbf{\ta{ஆசை}}? (\textbf{A wish, a longing.}) \textquotedblleft{}My goal\textquotedblright{}? (\textbf{\ta{என்} \ta{இலக்கு}}.)
\end{tcolorbox}
\section[\texorpdfstring{\ta{போக} \ta{ஆசை} (pōka ācai)}{போக ஆசை (pōka ācai)}]{\ta{போக} \ta{ஆசை} (pōka ācai) — what I hope for}

\label{lesson:TA-C290-poka-acai}



\begin{quote}
Say \textquotedblleft{}I long to swim in the sea.\textquotedblright{} (\emph{eṉakkuk kaṭalil nīnta ācai}.) Then \textquotedblleft{}My goal is to become a doctor.\textquotedblright{} (\emph{eṉ ilakku maruttuvar āvatu}.) Then \textquotedblleft{}I must go.\textquotedblright{} (\emph{nāṉ pōka vēṇṭum}.)
\end{quote}

\subsection*{You'll want to know: a wish, with the to-form}
\begin{quote}
\textbf{\ta{எனக்கு} \ta{ஒரு} \ta{நாள்} \ta{விமானத்தில்} \ta{பயணம்} \ta{செய்ய} \ta{ஆசை}.}
\end{quote}

\begin{quote}
\emph{eṉakku oru nāḷ vimāṉattil payaṇam ceyya ācai.} — \textquotedblleft{}I long to travel by plane one day.\textquotedblright{}
\end{quote}

\textbf{\ta{விமானத்தில்}} (\emph{vimāṉattil}, \textquotedblleft{}by plane\textquotedblright{}); \textbf{\ta{பயணம்} \ta{செய்}} (\emph{payaṇam cey}, \textquotedblleft{}to travel\textquotedblright{}), \textquotedblleft{}do a journey\textquotedblright{}. The to-form that came before \textbf{\ta{வேண்டும்}} comes before \textbf{\ta{ஆசை}}: \textbf{\ta{செய்ய} \ta{ஆசை}}.

\begin{quote}
\textbf{\ta{நான்} \ta{தமிழ்} \ta{நன்றாகப்} \ta{பேச} \ta{விரும்புகிறேன்}.}
\end{quote}

\begin{quote}
\emph{nāṉ tamiḻ naṉṟākap pēca virumpukiṟēṉ.} — \textquotedblleft{}I want to speak Tamil well.\textquotedblright{}
\end{quote}

\textbf{\ta{விரும்பு}}, \textquotedblleft{}to wish, to like\textquotedblright{}, takes a plain subject, \textbf{\ta{நான்}}, and the same to-form. It is plainer than \textbf{\ta{ஆசை}}, and fits a letter or a form as well as talk.

\subsection*{You'll want to know: what you want to become, \ta{ஆக}}
\begin{quote}
\textbf{\ta{நான்} \ta{ஆசிரியர்} \ta{ஆக} \ta{விரும்புகிறேன்}.}
\end{quote}

\begin{quote}
\emph{nāṉ āciriyar āka virumpukiṟēṉ.} — \textquotedblleft{}I want to become a teacher.\textquotedblright{}
\end{quote}

\textbf{\ta{ஆக}} (\emph{āka}, \textquotedblleft{}to become\textquotedblright{}) is the to-form of \textbf{\ta{ஆகு}}. The thing you will become stands bare in front of it: \textbf{\ta{ஆசிரியர்} \ta{ஆக}}, \textbf{\ta{மருத்துவர்} \ta{ஆக}}.

Three ways to say the same hope: \textbf{\ta{எனக்கு} \ta{ஆசிரியர்} \ta{ஆக} \ta{ஆசை}}; \textbf{\ta{நான்} \ta{ஆசிரியர்} \ta{ஆக} \ta{விரும்புகிறேன்}}; \textbf{\ta{என்} \ta{இலக்கு} \ta{ஆசிரியர்} \ta{ஆவது}}. And the biggest of them: \textbf{\ta{என்} \ta{கனவு}} (\emph{eṉ kaṉavu}, \textquotedblleft{}my dream\textquotedblright{}).

\subsection*{Your turn}
Say these aloud:

\begin{itemize}
\raggedright
  \item \textquotedblleft{}I long to travel by plane one day\textquotedblright{} — \emph{eṉakku oru nāḷ vimāṉattil payaṇam ceyya ācai}
  \item \textquotedblleft{}I want to speak Tamil well\textquotedblright{} — \emph{nāṉ tamiḻ naṉṟākap pēca virumpukiṟēṉ}
  \item \textquotedblleft{}I want to become a teacher\textquotedblright{} — \emph{nāṉ āciriyar āka virumpukiṟēṉ}
  \item one hope of your own, with \ta{ஆசை}, and one goal, with \ta{இலக்கு}
\end{itemize}

\begin{tcolorbox}[breakable,colback=teal!4,colframe=teal!35!black,title={Before you move on}]
\textquotedblleft{}I long to go\textquotedblright{}? (\textbf{\ta{எனக்குப்} \ta{போக} \ta{ஆசை}}.) \textquotedblleft{}I want to speak\textquotedblright{}? (\textbf{\ta{நான்} \ta{பேச} \ta{விரும்புகிறேன்}}.) \textquotedblleft{}To become\textquotedblright{}? (\textbf{\ta{ஆக}}.) Which one takes \textbf{-\ta{க்கு}}, and which one takes \textbf{\ta{நான்}}? (\textbf{\ta{ஆசை}} takes \textbf{-\ta{க்கு}}; \textbf{\ta{விரும்புகிறேன்}} takes \textbf{\ta{நான்}}.)
\end{tcolorbox}
\section[(review)]{Review — feelings and hopes}

\label{lesson:TA-R290-feelings-and-hopes}



\begin{quote}
Say \textquotedblleft{}a wish\textquotedblright{} and \textquotedblleft{}a goal\textquotedblright{}. Then read \textbf{\ta{ஈரம்}} aloud. (\emph{īram}.)

\begin{itemize}
\raggedright
  \item \emph{Write it:} \textbf{\ta{ஆசை}} and \textbf{\ta{இலக்கு}}, once each
\end{itemize}
\end{quote}

\subsection*{Your turn}
Say these aloud:

\begin{itemize}
\raggedright
  \item I missed the bus. I got angry.
  \item When my friend came, I was happy.
  \item I long to swim in the sea one day.
  \item My goal is to speak Tamil well.
  \item one feeling from last week, and why you felt it
\end{itemize}

\begin{tcolorbox}[breakable,colback=teal!4,colframe=teal!35!black,title={Before you move on}]
\textquotedblleft{}I got angry\textquotedblright{}? (\textbf{\ta{எனக்குக்} \ta{கோபம்} \ta{வந்தது}}.) \textquotedblleft{}I was happy\textquotedblright{}? (\textbf{\ta{எனக்குச்} \ta{சந்தோஷமாக} \ta{இருந்தது}}.) \textquotedblleft{}I long to go\textquotedblright{}? (\textbf{\ta{எனக்குப்} \ta{போக} \ta{ஆசை}}.) \textquotedblleft{}To become a doctor\textquotedblright{}? (\textbf{\ta{மருத்துவர்} \ta{ஆக}}.)
\end{tcolorbox}
